\documentclass{article}
\usepackage{amsmath,hyperref}
\title{Support handoff under two changes}
\date{}
\begin{document}
\maketitle
This primary civil warehouse case extends the reached \href{https://github.com/eternal-november/eos-research-programmes/pull/346}{connected-handoff pilot}. It is not C11 acceptance.

\paragraph{Plant.} At tick $-2$ loads 0--3 are each placed on their native support. At $-1$ previously scheduled rack settling removes load 2, and in the short case also load 1. Load 3 stays put. A local authenticated warehouse order known at $-1$ schedules authorization of root load 0 for tick 0, when the request occurs in both policies. The warehouse is owed valid relief of root load 0 by tick 5 while all live loads stay supported. Each load may move one position to its successor. No bypass, additional routes, corrupt actors, clock drift or independent loss or delay of the vacancy signal is admitted. Each pending grant keeps both endpoints reserved and the old physical support in place until actual acceptance. The atomic adjacent-endpoint guard makes the corresponding vacancy visible at acceptance; release and visibility are distinct effects that coincide in this plant.

\paragraph{Event phase and disturbance.} Process due accepts first, then channel events, then new starts. The tick-0 first-hop queue onset delays packets due then if it wins the race, and delays the tick-0 issue if present; retiming never occurs after delivery. The packet sent at $-1$ was enqueued for 0 without knowledge of this edit. The second onset at 0 lengthens terminal dwell into support 2 from one to two ticks; no start races with that onset. An authorized receiver forwards only on receipt while holding its native load. A start requires the named permit to be active, the source load and free target to be present, neither endpoint to be reserved, and the clock to be before deadline. These are receiver-local facts from the adjacent-endpoint interlock; no controller consults the arm or future dwell. Due acceptances remain in effect after deadline; no fresh start is permitted then. A missed deadline is not an authorized refusal.

\paragraph{Results.} In nominal, hop-only, dwell-only, and combined long arms, service times are 4, 5, 5, and missed. In the combined arm support 1 receives at 2 and issues forward; support 2 receives and begins at 3; physical acceptance at 5 frees support 1 too late to start root relief. These two interfaces carry actual propagated effects under the installed controller. At the short reached cut, the prior rack settlement frees the same middle support: service ticks are 2, 3, 2, 3 with identical topology and edits. This separates the retained-work effect, not an isolated mediator intervention at fixed load; turning off dwell would merely remove the initiating edit. Deadline 6 is a nearby nonfailure and with deadline 4 neither long single edit meets service. Both are sensitivity controls, not changes to the tick-5 promise.

\paragraph{Repair and lower bound.} With issue at or after 0, receipt before forwarding, unchanged dwell and exclusive adjacent support, even a clairvoyant scheduler cannot achieve service: first/second receipts are at least 2/3, middle release at least 5, and root requires one further tick. A permissible repair on the declared workload sends the committed permit at $-1$ with activation at 0. The long times are 3, 4, 4, 5 with onset first and 3, 3, 4, 4 with delivery first. Both short cases also serve. This changes only the local root sender, with one tick of preparation and earlier use of the same channel, no extra message or support. Authorization before a surprise order is not established. Unchanged receivers receive changed authority timing and physical availability; both must be checked.

\paragraph{Reference and regression.} One-tick grants in the original plant never persist into a later start phase, so per-tick claims suffice there. Dwell-only extension duplicates a terminal grant and produces a stale-acceptance assertion at 5. Persistent endpoint exclusion is common to all arms: removing due grants and then adding only disjoint endpoints preserves disjointness inductively. Existing support persists until actual acceptance, never merely until deadline. The executable reference independently enumerates allowable start ticks (or terminal/root pairs), testing the lower bounds on authority and physical release, not scheduling through the event loop. It matches the simulator across 96 combinations, including the two orders and deadlines 4--6. Nominal outcomes match upstream service and final stock on the six distinct one-tick controls.

\paragraph{Proposed proof boundary, not frozen.} Use the pinned \href{https://github.com/eternal-november/eos-research-programmes/blob/e1348964f73704212d2d8b48d5982a5ed600a2ff/projects/contextual-composition/feedback-and-concurrency/observable-covers/Residual_Profiles_and_Admission_Reflection.tex}{Residual Profiles and Admission Reflection}, theorem Existence and construction of common profiles. Receiver projection preserves actual order receipt, identity and activation, clock and deadline, adjacent stocks and reservations, start and actual acceptance. The latter are different events. A deadline bit loses differences potentially consumed by unchanged timed rules. Candidate native witnesses are channel/authority, physical support/reservations and order/clock histories. Actual native maps, primitive-test reflection on full declared domains, residual class coverage and a sequential invariant are open. Enumerated successful traces do not imply any of them; independently selected positive and rejected transfers remain owed. Do not give the receiver new inputs to obtain coverage. The modular direct proof uses the same cached local guard lemma and rechecks authority, timed resource availability and endpoint exclusion. No proof saving is claimed.
\end{document}
